%%
%% AgentTelemetry: Towards Unified Observability for Autonomous AI Agent Systems
%% Position Paper — Target venue: ICSE 2027 SEIP / ASE 2027 Workshop
%%

\documentclass[sigconf,nonacm]{acmart}

%% ---------- suppress acmart warnings for non-ACM use ----------
\setcopyright{none}
\settopmatter{printacmref=false}
\makeatletter
\@ifundefined{footnotetextcopyrightpermission}{}{%
  \renewcommand\footnotetextcopyrightpermission[1]{}%
}
\makeatother
\pagestyle{plain}

%% ---------- packages ----------
\usepackage{listings}
\usepackage{xcolor}
\usepackage{booktabs}
\usepackage{enumitem}
\usepackage{url}

%% ---------- listings style ----------
\lstset{
  language=Python,
  basicstyle=\ttfamily\small,
  keywordstyle=\color{blue!70!black}\bfseries,
  commentstyle=\color{green!50!black},
  stringstyle=\color{red!60!black},
  showstringspaces=false,
  breaklines=true,
  breakatwhitespace=true,
  frame=single,
  framesep=3pt,
  xleftmargin=6pt,
  xrightmargin=4pt,
  numbers=left,
  numberstyle=\tiny\color{gray},
  numbersep=5pt,
  aboveskip=8pt,
  belowskip=6pt,
  columns=flexible,
}

%% ---------- metadata ----------
\title{AgentTelemetry: Towards Unified Observability for Autonomous AI Agent Systems}

\author{Krishna Chaitanya Balusu}
\affiliation{%
  \institution{Independent Researcher}
  \country{United States}
}
\email{krishna.balusu@outlook.com}

\begin{abstract}
Autonomous AI agents built on large language models are rapidly transitioning
from research prototypes to production deployments, yet the observability
infrastructure required to monitor, debug, and govern these systems remains
fundamentally inadequate.  Existing observability standards such as
OpenTelemetry were designed for request-response microservice architectures and
lack semantic primitives for agent-specific phenomena including multi-step
reasoning chains, tool invocation sequences, inter-agent delegation, and
cost accumulation across LLM calls.  We present \textsc{AgentTelemetry}, an
open-source, framework-agnostic observability library that introduces (1) a
telemetry data model with nine agent-specific span kinds and semantic
attribute conventions covering agent identity, LLM usage, tool execution, and
inter-agent communication; (2) a W3C-compatible context propagation mechanism
for end-to-end tracing across multi-agent topologies; (3) a privacy-first
design in which prompt and completion capture is strictly opt-in; and (4) a
pluggable instrumentor--exporter architecture supporting automatic telemetry
collection from LangChain, CrewAI, AutoGen, and DSPy with export to console,
JSON, and any OpenTelemetry Protocol (OTLP) backend.  We describe the system
design, contrast it with existing approaches in the literature, and outline a
preliminary evaluation plan measuring instrumentation overhead and
debugging effectiveness.
\end{abstract}

\begin{CCSXML}
<ccs2012>
<concept>
<concept_id>10011007.10011074.10011099</concept_id>
<concept_desc>Software and its engineering~Software verification and validation</concept_desc>
<concept_significance>500</concept_significance>
</concept>
<concept>
<concept_id>10011007.10011074.10011111</concept_id>
<concept_desc>Software and its engineering~Software post-development issues</concept_desc>
<concept_significance>300</concept_significance>
</concept>
</ccs2012>
\end{CCSXML}

\ccsdesc[500]{Software and its engineering~Software verification and validation}
\ccsdesc[300]{Software and its engineering~Software post-development issues}

\keywords{AI agents, observability, telemetry, distributed tracing, LLM monitoring, OpenTelemetry, multi-agent systems}

\begin{document}

\maketitle

%% ====================================================================
%% 1. INTRODUCTION
%% ====================================================================
\section{Introduction}
\label{sec:introduction}

Large language model (LLM) based autonomous agents represent a paradigm shift
in software architecture.  Unlike conventional services that receive a request,
execute deterministic logic, and return a response, an LLM agent pursues a
goal through an iterative cycle of reasoning, planning, tool use, and
self-reflection---often delegating sub-tasks to other agents in a multi-agent
topology~\cite{agentbench,agentboard}.  Frameworks such as LangChain,
CrewAI, AutoGen, and DSPy have lowered the barrier to building such agents,
and organizations are deploying them for customer support, code generation,
data analysis, and autonomous operations~\cite{aiops_survey}.

This proliferation has exposed a critical gap: \emph{there is no
standardized observability infrastructure for AI agent systems}.
Production teams cannot reliably answer fundamental questions about their
agents: Why did the agent choose this tool over that one?  Where in the
reasoning chain did the hallucination originate?  How much did this task
cost across all LLM calls?  Which agent in the multi-agent pipeline is the
bottleneck?

The observability industry has converged on OpenTelemetry
(OTel)~\cite{otel} as the standard for distributed tracing, metrics, and
logging in microservice architectures.  However, OTel's data model was
designed for request-response service interactions: its \texttt{SpanKind}
enum (\texttt{SERVER}, \texttt{CLIENT}, \texttt{PRODUCER}, \texttt{CONSUMER},
\texttt{INTERNAL}) has no representation for reasoning steps, LLM calls,
tool invocations, or inter-agent messages.  Its semantic conventions
define attributes for HTTP, gRPC, and database operations but not for
token counts, model selection, prompt content, or cost estimation.

Recent research has begun to address the observability of AI systems from
various angles.  AgentSight~\cite{agentsight} uses eBPF to bridge the
semantic gap between high-level agent intent and kernel-level system
actions, providing security-focused monitoring.  NeMo
Guardrails~\cite{nemo_guardrails} offers programmable runtime safety
rails that intercept agent inputs and outputs.  DSPy
Assertions~\cite{dspy_assertions} embed computational constraints into
LLM pipelines to trigger self-refinement on violations.
AgentBoard~\cite{agentboard} introduces fine-grained progress tracking
across sub-goals in multi-turn agent tasks.  Agent-as-a-Judge~\cite{agent_as_judge}
proposes using one agent to evaluate another's trajectory.
These contributions address important facets---security monitoring, safety
enforcement, constraint checking, evaluation, and trajectory
assessment---but none provides a \emph{general-purpose telemetry framework}
that unifies tracing, metrics, and context propagation specifically for
agent workloads.

In this paper, we present \textsc{AgentTelemetry}, an open-source library
that fills this gap.  Our contributions are:
\begin{enumerate}[leftmargin=*,nosep]
    \item A \textbf{telemetry data model} with nine agent-specific span
        kinds, semantic attribute conventions for agent identity, LLM
        usage, tool execution, and inter-agent communication, and
        built-in cost estimation.
    \item A \textbf{context propagation protocol} compatible with the W3C
        Trace Context standard, extended to carry agent-specific metadata
        across multi-agent boundaries.
    \item A \textbf{privacy-first design} in which prompt and completion
        content is never captured unless the operator explicitly opts in.
    \item A \textbf{pluggable instrumentor--exporter architecture} that
        decouples framework-specific telemetry collection from backend
        export, with concrete implementations for four major agent
        frameworks and three export targets including OTLP.
\end{enumerate}

%% ====================================================================
%% 2. MOTIVATION AND GAP ANALYSIS
%% ====================================================================
\section{Motivation and Gap Analysis}
\label{sec:motivation}

\subsection{The Agent Observability Problem}
An LLM agent's execution produces a deeply nested, non-deterministic trace
of cognitive operations.  Consider a multi-agent research assistant: a
\emph{planner} agent decomposes a user question into sub-tasks, a
\emph{researcher} agent performs web searches and retrieves documents, and
a \emph{writer} agent synthesizes findings.  Each agent may invoke
multiple LLM calls with different models and temperature settings, use
various tools (web search, code execution, file I/O), and exchange messages
with peer agents.  A single user request can produce dozens of spans
across three or more agents, accumulating thousands of tokens and
non-trivial cost.

Debugging a failure in this pipeline---for example, the writer
producing a factually incorrect summary---requires tracing backward
through the writer's LLM calls to the researcher's retrieval results to
the planner's task decomposition.  Without unified tracing across agent
boundaries, operators resort to ad-hoc logging, which is inconsistent
across frameworks, omits critical context, and scales poorly.

\subsection{Why OpenTelemetry Is Necessary but Insufficient}
OpenTelemetry provides a robust foundation for distributed tracing: trace
and span identifiers, context propagation, exporters, and a rich
ecosystem of backends.  \textsc{AgentTelemetry} is deliberately designed
to be OTel-compatible; its spans map one-to-one to OTel spans when
exported via the OTLP exporter.

However, raw OTel lacks the \emph{semantic layer} needed for agent
observability.  Table~\ref{tab:gap} summarizes the conceptual mapping
and the gaps.  OTel's \texttt{SpanKind} has no equivalent for reasoning,
planning, reflection, or retrieval steps.  Its semantic conventions lack
attributes for token counts, model identity, cost, and agent role.
Multi-agent context propagation requires carrying agent identity and
delegation metadata that the W3C \texttt{traceparent} header does not encode.

\begin{table}[t]
\centering
\small
\caption{Mapping of agent concepts to observability primitives.  Bold entries indicate novel contributions of \textsc{AgentTelemetry}.}
\label{tab:gap}
\begin{tabular}{@{}lll@{}}
\toprule
\textbf{Agent Concept} & \textbf{OTel Equivalent} & \textbf{AgentTelemetry} \\
\midrule
Agent            & Service         & \texttt{agent.name, .role} \\
Task execution   & Trace           & \textbf{AgentTrace} \\
Reasoning step   & (none)          & \textbf{REASONING span} \\
LLM call         & (none)          & \textbf{LLM\_CALL span} \\
Tool invocation  & (none)          & \textbf{TOOL\_CALL span} \\
Planning         & (none)          & \textbf{PLANNING span} \\
Self-reflection  & (none)          & \textbf{REFLECTION span} \\
Retrieval (RAG)  & (none)          & \textbf{RETRIEVAL span} \\
Agent message    & (none)          & \textbf{AGENT\_COMM span} \\
Human-in-loop    & (none)          & \textbf{HUMAN\_INPUT span} \\
Token count      & (none)          & \texttt{llm.input/output\_tokens} \\
Cost             & (none)          & \texttt{llm.cost\_usd} \\
Prompt content   & (none)          & \texttt{llm.prompt} (opt-in) \\
\bottomrule
\end{tabular}
\end{table}

\subsection{Limitations of Existing Approaches}
AgentSight~\cite{agentsight} operates at the kernel level via eBPF,
monitoring TLS-encrypted LLM traffic and system calls.  This provides
valuable security insights but requires Linux-specific kernel support,
cannot capture framework-level semantics (which tool was called, what
the agent's reasoning was), and introduces a dependency on
kernel-level instrumentation that many cloud deployments restrict.

NeMo Guardrails~\cite{nemo_guardrails} and DSPy
Assertions~\cite{dspy_assertions} focus on \emph{controlling} agent
behavior through runtime constraints rather than \emph{observing} it.
They provide valuable safety enforcement but do not produce structured
telemetry suitable for dashboards, alerting, or post-hoc analysis.

Evaluation frameworks such as AgentBoard~\cite{agentboard},
AgentBench~\cite{agentbench}, and Agent-as-a-Judge~\cite{agent_as_judge}
are designed for offline assessment of agent capabilities, not for
continuous production monitoring.

Monitoring and observability research for ML systems~\cite{ml_monitoring}
has documented significant gaps in practice but has not proposed a
concrete telemetry standard.  Safety evaluation
frameworks~\cite{openagentsafety, toolemu} identify failure modes but do
not address operational observability.

\textsc{AgentTelemetry} occupies a distinct position: it is a
\emph{production-oriented, framework-agnostic telemetry library} that
provides structured observability primitives specifically designed for
agent workloads, bridging the gap between existing distributed tracing
standards and the unique requirements of autonomous AI systems.

%% ====================================================================
%% 3. DESIGN
%% ====================================================================
\section{System Design}
\label{sec:design}

\textsc{AgentTelemetry} is organized into three layers: a \emph{core
telemetry model}, a \emph{framework instrumentor layer}, and an
\emph{exporter layer}.  Figure~\ref{fig:architecture} describes the
conceptual architecture.

\begin{figure}[t]
\centering
\fbox{\parbox{0.92\columnwidth}{\small\ttfamily
~~Agent Framework (LangChain, CrewAI, ...)\\
~~~~~~~~~~\textnormal{$\downarrow$}~~auto-instrumentation \\
~~Instrumentor Layer (BaseInstrumentor)\\
~~~~~~~~~~\textnormal{$\downarrow$}~~AgentSpan, AgentEvent \\
~~Core Telemetry (AgentTracer, AgentMetrics)\\
~~~~~~~~~~\textnormal{$\downarrow$}~~export\_span() \\
~~Exporter Layer (Console | JSON | OTLP)\\
~~~~~~~~~~\textnormal{$\downarrow$}\\
~~Backend (Grafana Tempo, Jaeger, Datadog, ...)
}}
\caption{Layered architecture of \textsc{AgentTelemetry}.  Instrumentors are
framework-specific; the core and exporter layers are framework-agnostic.}
\label{fig:architecture}
\end{figure}

\subsection{Telemetry Data Model}
\label{sec:datamodel}

\paragraph{Span Kinds.}
The \texttt{AgentSpanKind} enumeration defines nine span types:
\texttt{TASK} (root span for a complete task), \texttt{REASONING} (chain-of-thought),
\texttt{LLM\_CALL} (language model invocation), \texttt{TOOL\_CALL}
(function/tool execution), \texttt{PLANNING} (task decomposition),
\texttt{REFLECTION} (self-critique), \texttt{RETRIEVAL} (RAG lookup),
\texttt{AGENT\_COMM} (inter-agent message), and \texttt{HUMAN\_INPUT}
(human-in-the-loop interaction).  Each kind maps to a distinct phase in
the agent execution lifecycle and enables kind-specific dashboards and
alerts.

\paragraph{Semantic Attribute Conventions.}
We define three attribute namespaces.  The \texttt{agent.*} namespace
captures agent identity: \texttt{agent.name}, \texttt{agent.framework},
\texttt{agent.framework.version}, \texttt{agent.role}, and
\texttt{agent.task}.  The \texttt{llm.*} namespace captures LLM call
metadata: \texttt{llm.model}, \texttt{llm.provider},
\texttt{llm.input\_tokens}, \texttt{llm.output\_tokens},
\texttt{llm.total\_tokens}, \texttt{llm.temperature},
\texttt{llm.cost\_usd}, and \texttt{llm.latency\_ms}.  The
\texttt{tool.*} namespace captures tool execution details:
\texttt{tool.name}, \texttt{tool.description}, \texttt{tool.input},
\texttt{tool.output}, \texttt{tool.success}, \texttt{tool.error}, and
\texttt{tool.latency\_ms}.  Additional attributes for inter-agent
interactions record source and target agent identifiers and interaction
type.

\paragraph{Cost Estimation.}
Each \texttt{LLM\_CALL} span automatically estimates cost in USD based
on the model name and token counts.  The tracer maintains a lookup table
of per-million-token prices for major models (GPT-4o, GPT-4-turbo,
Claude 3.5 Sonnet, Claude Opus 4, etc.)\ and computes cost at span
completion.  This enables aggregate cost dashboards without requiring
external billing API integration.

\paragraph{Events.}
Spans can contain timestamped \texttt{AgentEvent} instances categorized
by \texttt{EventType}: \texttt{LLM\_START}, \texttt{LLM\_END},
\texttt{TOOL\_START}, \texttt{TOOL\_END}, \texttt{GUARDRAIL\_TRIGGERED},
\texttt{COST\_THRESHOLD}, \texttt{ERROR}, \texttt{WARNING}, and others.
Events provide fine-grained intra-span signals---for example, recording when
a guardrail fires during an LLM call or when cumulative cost exceeds a
configured threshold.

\subsection{Context Propagation for Multi-Agent Systems}
\label{sec:context}

Multi-agent architectures require trace context to flow across agent
boundaries so that spans from different agents share a single trace
identity.  \textsc{AgentTelemetry} provides the \texttt{AgentContext}
class, which follows the W3C Trace Context format and extends it with
agent-specific metadata.

The context is serialized to a carrier dictionary containing three
entries: (1)~a \texttt{traceparent} header in W3C format
(\texttt{00-<trace\_id>-<parent\_span\_id>-01}); (2)~an
\texttt{agentstate} header carrying the source agent's name; and
(3)~an optional \texttt{baggage} header for key-value metadata
propagation.  The carrier can be transmitted via HTTP headers, message
queue metadata, or any string-based transport.

On the receiving side, \texttt{AgentContext.from\_carrier()} deserializes
the context, and the receiving agent's tracer attaches it so that
subsequent spans inherit the trace~ID and parent span reference.  This
produces a unified trace tree that spans multiple agents, enabling
operators to visualize the complete execution flow of a multi-agent
task in a single trace viewer.

\subsection{Privacy-First Design}
\label{sec:privacy}

LLM prompts and completions frequently contain sensitive information:
user queries, personal data, proprietary documents, or internal system
instructions.  Capturing this content in telemetry creates significant
privacy and compliance risk.

\textsc{AgentTelemetry} adopts a privacy-first posture: the
\texttt{llm.prompt} and \texttt{llm.completion} attributes are defined in
the semantic conventions but are \emph{never populated} unless the operator
explicitly sets \texttt{capture\_content=True} when initializing the
tracer.  This opt-in gate applies at both the tracer and instrumentor
levels, providing defense in depth.

When content capture is disabled, the telemetry still records all
structural information---span kind, duration, token counts, cost, model
identity, tool names, success/failure status---providing comprehensive
observability without exposing sensitive content.  This design aligns
with emerging regulatory requirements around AI system transparency and
enables deployment in regulated industries (healthcare, finance) where
data minimization is mandatory.

\subsection{Metrics Collection}
\label{sec:metrics}

Complementing per-span trace data, the \texttt{AgentMetrics} class
provides thread-safe aggregate metrics through three instrument types:
counters (e.g., \texttt{agent.llm.call.count},
\texttt{agent.cost.total\_usd}), histograms (e.g.,
\texttt{agent.llm.latency\_ms}, \texttt{agent.tool.latency\_ms}), and
gauges.  Each metric point carries a label set including
\texttt{agent.name} and optional dimensions (model name, tool name),
enabling multi-dimensional analysis.  The \texttt{summary()} method
computes descriptive statistics (min, max, mean, p50, p95, p99) over
histogram data, providing immediate insight into latency distributions
without requiring an external analytics backend.

%% ====================================================================
%% 4. IMPLEMENTATION
%% ====================================================================
\section{Implementation}
\label{sec:implementation}

\subsection{Framework-Agnostic Instrumentor Architecture}
\label{sec:instrumentor}

The instrumentor layer is designed around a template method pattern.
\texttt{BaseInstrumentor} is an abstract class that provides:

\begin{enumerate}[leftmargin=*,nosep]
    \item A default \texttt{AgentTracer} and \texttt{AgentMetrics}
        instance, configured with the framework name.
    \item Convenience methods \texttt{\_record\_llm\_metrics()} and
        \texttt{\_record\_tool\_metrics()} that atomically update
        counters and histograms for LLM and tool calls.
    \item Abstract methods \texttt{instrument()} and
        \texttt{uninstrument()} that subclasses implement with
        framework-specific monkey-patching or callback registration.
    \item A \texttt{capture\_content} flag that propagates the
        privacy setting to the underlying tracer.
\end{enumerate}

Listing~\ref{lst:instrumentor_usage} shows the usage pattern.  The
instrumentor applies patches transparently; user code interacts with
the agent framework as normal and telemetry is collected automatically.

\begin{lstlisting}[caption={Instrumentor usage: automatic telemetry
  collection requires no changes to user code.},
  label={lst:instrumentor_usage},float=t]
from agenttelemetry.instrumentors.langchain \
    import LangChainInstrumentor
from agenttelemetry.exporters.otlp \
    import OTLPExporter

# One-time setup
inst = LangChainInstrumentor()
inst.tracer.add_exporter(
    OTLPExporter(endpoint="otel-collector:4317")
)
inst.instrument()

# Use LangChain normally -- spans are
# captured and exported automatically
agent = create_langchain_agent(...)
result = agent.invoke("Summarize this paper")

# Teardown
inst.uninstrument()
\end{lstlisting}

\subsection{Supported Frameworks}
\label{sec:frameworks}

We target four major agent frameworks that collectively represent the
majority of production agent deployments:

\begin{description}[leftmargin=*,nosep,style=unboxed]
    \item[LangChain/LangGraph.]  Callback-based instrumentation via
        LangChain's \texttt{BaseCallbackHandler}.  Captures chain
        invocations, LLM calls, tool calls, and retriever operations
        as nested spans.
    \item[CrewAI.]  Patches \texttt{Crew.kickoff()} and agent task
        execution methods to capture multi-agent role assignments,
        delegation events, and per-agent cost breakdowns.
    \item[AutoGen.]  Hooks into the \texttt{ConversableAgent} message
        handling and \texttt{GroupChat} orchestration to trace
        inter-agent conversations and tool use within chat turns.
    \item[DSPy.]  Wraps DSPy module \texttt{forward()} methods and
        language model calls to capture the compilation--optimization
        pipeline alongside runtime execution traces.
\end{description}

Each instrumentor produces spans with the same semantic conventions,
ensuring that dashboards and alerts work identically regardless of the
underlying framework.

\subsection{Exporter Ecosystem}
\label{sec:exporters}

Three exporters are provided:

\begin{description}[leftmargin=*,nosep,style=unboxed]
    \item[ConsoleExporter.]  Formats spans as single-line summaries
        showing status, span kind, name, duration, token counts, cost,
        and tool success.  A verbose mode emits full JSON.  Designed
        for development and debugging.
    \item[JSONFileExporter.]  Writes spans as JSON Lines to a file,
        enabling offline analysis, integration with log management
        tools, and archival.
    \item[OTLPExporter.]  Converts \texttt{AgentSpan} instances to
        OpenTelemetry spans and exports them via gRPC to any
        OTLP-compatible backend.  Agent-specific attributes are
        preserved as OTel span attributes; agent events map to OTel
        events.  This enables visualization in Grafana Tempo, Jaeger,
        Zipkin, Datadog, Honeycomb, and other platforms with no
        additional configuration.
\end{description}

The exporter interface is minimal---a single \texttt{export\_span(span)}
method---making it straightforward to implement custom exporters for
proprietary backends or specialized analytics pipelines.

\subsection{Tracing API}
\label{sec:tracing_api}

The \texttt{AgentTracer} class is the primary user-facing API.  It
provides context-manager methods (\texttt{start\_task},
\texttt{start\_llm\_call}, \texttt{start\_tool\_call},
\texttt{start\_reasoning}, \texttt{start\_planning},
\texttt{start\_retrieval}, \texttt{start\_agent\_comm}) that create
properly nested spans with automatic parent--child linking, timing,
error capture, and cost estimation.  Spans are automatically exported
to all registered exporters upon completion.  Listing~\ref{lst:tracing}
illustrates direct API usage for custom agent implementations.

\begin{lstlisting}[caption={Direct tracing API for custom agents.
  Spans nest automatically via context managers.},
  label={lst:tracing},float=t]
tracer = AgentTracer(
    agent_name="researcher",
    framework="custom",
    capture_content=False
)
tracer.add_exporter(ConsoleExporter())

with tracer.start_task("Analyze dataset") as t:
  with tracer.start_planning("decompose") as p:
    pass  # planning logic
  with tracer.start_llm_call(model="gpt-4o") as l:
    l.set_attribute("llm.input_tokens", 1200)
    l.set_attribute("llm.output_tokens", 450)
    # cost auto-calculated on span end
  with tracer.start_tool_call("sql_query") as tc:
    tc.set_attribute("tool.input", "SELECT ...")
    tc.set_attribute("tool.success", True)
\end{lstlisting}

%% ====================================================================
%% 5. PRELIMINARY EVALUATION
%% ====================================================================
\section{Preliminary Evaluation}
\label{sec:evaluation}

As this is a position paper presenting the design and initial implementation
of \textsc{AgentTelemetry}, we outline our planned evaluation rather than
reporting completed experiments.

\subsection{Instrumentation Overhead}
\label{sec:overhead}

We plan to measure the runtime overhead of \textsc{AgentTelemetry}
instrumentation across the four supported frameworks.  The experimental
design is as follows:

\begin{itemize}[leftmargin=*,nosep]
    \item \textbf{Workloads:}  Standard benchmark tasks from
        AgentBench~\cite{agentbench} (code generation, web browsing,
        database querying) executed with and without instrumentation.
    \item \textbf{Metrics:}  Wall-clock time per task, peak memory usage,
        and CPU utilization delta.
    \item \textbf{Exporters:}  Overhead will be measured separately for
        each exporter (console, JSON, OTLP) to isolate I/O costs.
    \item \textbf{Hypothesis:}  We expect instrumentation overhead to be
        dominated by LLM API latency (typically 500ms--5s per call),
        making the microsecond-scale cost of span creation and
        attribute setting negligible in practice ($<$0.1\% overhead).
\end{itemize}

\subsection{Debugging Effectiveness Case Study}
\label{sec:case_study}

We plan a controlled case study in which participants debug agent
failures with and without \textsc{AgentTelemetry}.  The study will use
a multi-agent research pipeline (planner $\rightarrow$ researcher
$\rightarrow$ writer) with injected faults: (1)~a hallucinated retrieval
result that propagates to the final output, (2)~a tool timeout that
causes the agent to retry with degraded context, and (3)~a cost overrun
caused by excessive LLM calls in a reasoning loop.

For each fault, we will measure time-to-diagnosis and diagnostic accuracy
with three conditions: (a)~framework-default logging only, (b)~\textsc{AgentTelemetry}
with console export, and (c)~\textsc{AgentTelemetry} with OTLP export
to a Grafana dashboard.  We hypothesize that structured telemetry with
agent-specific span kinds will significantly reduce diagnostic time,
particularly for cross-agent fault propagation scenarios.

%% ====================================================================
%% 6. RELATED WORK
%% ====================================================================
\section{Related Work}
\label{sec:related}

\paragraph{System-Level Agent Monitoring.}
AgentSight~\cite{agentsight} is the most closely related work.  It uses
eBPF to monitor TLS-encrypted LLM traffic and correlate kernel events
with agent intent.  While AgentSight provides deep system-level
visibility with minimal overhead, it is restricted to Linux, requires
kernel-level access often unavailable in containerized deployments,
and cannot capture application-level semantics such as agent role,
task decomposition, or framework-specific context.
\textsc{AgentTelemetry} operates at the application level, providing
richer semantic telemetry at the cost of requiring framework integration.
The two approaches are complementary: AgentSight for security and
system-level monitoring, \textsc{AgentTelemetry} for application-level
observability.

\paragraph{Runtime Safety and Constraints.}
NeMo Guardrails~\cite{nemo_guardrails} provides programmable rails that
intercept and validate agent behavior in real time.  DSPy
Assertions~\cite{dspy_assertions} embed computational constraints into
LLM pipelines, triggering self-refinement on violations.  Both systems
enforce behavioral policies but do not produce the structured,
exportable telemetry needed for dashboards, trend analysis, or SLO
monitoring.  \textsc{AgentTelemetry}'s event system can record guardrail
activations (\texttt{GUARDRAIL\_TRIGGERED} events) within spans,
providing observability \emph{of} the guardrail system itself.

\paragraph{Agent Evaluation and Benchmarking.}
AgentBench~\cite{agentbench} evaluates LLMs as agents across eight
environments.  AgentBoard~\cite{agentboard} provides fine-grained
progress tracking across sub-goals.  Agent-as-a-Judge~\cite{agent_as_judge}
uses agents to evaluate peer agents' trajectories.  These frameworks
target offline evaluation and capability assessment, while
\textsc{AgentTelemetry} targets continuous production monitoring.
However, evaluation results from AgentBench can serve as benchmark
workloads for measuring instrumentation overhead (Section~\ref{sec:overhead}).

\paragraph{Agent Safety.}
ToolEmu~\cite{toolemu} emulates tool execution in a sandbox for
automated safety testing.  OpenAgentSafety~\cite{openagentsafety}
evaluates agent safety across 350+ multi-turn tasks and found unsafe
behavior in over 51\% of safety-vulnerable tasks.  These findings
underscore the need for production monitoring: safety evaluation during
development is necessary but not sufficient, and runtime observability
is essential for detecting unsafe behaviors that emerge only in
production contexts.

\paragraph{ML System Monitoring.}
Leest et al.~\cite{ml_monitoring} conducted focus groups documenting
current ML monitoring practices and gaps, finding that practitioners
lack standardized tools for ML-specific observability.  Their findings
motivate our work: \textsc{AgentTelemetry} addresses this gap for the
specific case of LLM-based agent systems.

\paragraph{AIOps and Distributed Tracing.}
The AIOps literature~\cite{aiops_survey, opsagent} has developed
sophisticated techniques for anomaly detection, root cause analysis,
and automated remediation in microservice architectures.  Systems like
TraceDiag~\cite{tracediag} and DeepTraLog~\cite{deeptralog} analyze
distributed traces for fault diagnosis.  \textsc{AgentTelemetry}'s OTLP
export capability enables these existing AIOps tools to be applied
directly to agent telemetry data, creating a bridge between the agent
and AIOps research communities.

\paragraph{Execution Tracing for Optimization.}
Trace~\cite{trace_opt} treats agent workflows analogously to
differentiable programs, tracing execution paths for automatic
optimization.  While Trace focuses on optimization through gradient-like
updates, \textsc{AgentTelemetry} focuses on observability through
structured telemetry export.  The recorded execution paths could
potentially be combined: \textsc{AgentTelemetry} traces feeding into
Trace-style optimizers.

%% ====================================================================
%% 7. CONCLUSION AND FUTURE WORK
%% ====================================================================
\section{Conclusion and Future Work}
\label{sec:conclusion}

We have presented \textsc{AgentTelemetry}, an open-source framework for
unified observability of autonomous AI agent systems.  By introducing
agent-specific span kinds, semantic attribute conventions, W3C-compatible
context propagation, and a privacy-first design, \textsc{AgentTelemetry}
bridges the gap between existing distributed tracing standards and the
unique requirements of LLM-based agent workloads.  Its pluggable
instrumentor--exporter architecture enables automatic telemetry collection
from major agent frameworks with export to any OTLP-compatible backend.

Several directions for future work are planned:

\begin{description}[leftmargin=*,nosep,style=unboxed]
    \item[Semantic Convention Standardization.]  We plan to propose
        \textsc{AgentTelemetry}'s semantic conventions as an OpenTelemetry
        Semantic Conventions Enhancement Proposal (OTEP), working toward
        an industry standard for agent telemetry.
    \item[Adaptive Sampling.]  Intelligent trace sampling that
        preferentially retains anomalous, high-cost, or failing traces
        while reducing storage volume for routine
        executions~\cite{sieve_sampling}.
    \item[Agent-Aware Dashboards.]  Pre-built Grafana dashboards that
        visualize agent-specific metrics: reasoning chain depth, tool
        usage patterns, cost accumulation curves, and multi-agent
        interaction graphs.
    \item[Streaming Telemetry.]  Real-time span streaming for live
        monitoring of long-running agent tasks, enabling operators to
        intervene before failures compound.
    \item[Safety Integration.]  First-class integration with NeMo
        Guardrails~\cite{nemo_guardrails} and DSPy
        Assertions~\cite{dspy_assertions} to record safety constraint
        violations as structured telemetry events within agent traces.
    \item[Causal Analysis.]  Applying causal inference techniques
        from the AIOps literature~\cite{mulan} to agent telemetry for
        automated root cause analysis of agent failures.
\end{description}

The growing deployment of autonomous AI agents in production environments
demands observability infrastructure that matches the complexity of these
systems.  \textsc{AgentTelemetry} provides a foundation for this
infrastructure, and we invite the research and practitioner communities to
contribute to its development.

%% ====================================================================
%% BIBLIOGRAPHY
%% ====================================================================

\begin{thebibliography}{20}

\bibitem{agentsight}
Y.~Zheng, Y.~Hu, T.~Yu, and A.~Quinn.
\newblock {AgentSight}: System-Level Observability for {AI} Agents Using {eBPF}.
\newblock {\em arXiv preprint arXiv:2508.02736}, 2025.

\bibitem{nemo_guardrails}
T.~Rebedea, R.~Dinu, M.~Sreedhar, C.~Parisien, and J.~Cohen.
\newblock {NeMo Guardrails}: A Toolkit for Controllable and Safe {LLM}
  Applications.
\newblock In {\em Proc.\ EMNLP (Demo Track)}, 2023.

\bibitem{dspy_assertions}
A.~Singhvi, M.~Shetty, S.~Tan, C.~Potts, K.~Sen, M.~Zaharia, and O.~Khattab.
\newblock {DSPy} Assertions: Computational Constraints for Self-Refining
  Language Model Pipelines.
\newblock In {\em Proc.\ COLM}, 2024.

\bibitem{agentboard}
C.~Ma, J.~Zhang, Z.~Zhu, C.~Yang, Y.~Yang, et~al.
\newblock {AgentBoard}: An Analytical Evaluation Board of Multi-Turn {LLM}
  Agents.
\newblock In {\em Proc.\ ACL}, 2024.

\bibitem{agent_as_judge}
M.~Zhuge, C.~Zhao, D.~Ashley, et~al.
\newblock Agent-as-a-Judge: Evaluate Agents with Agents.
\newblock {\em arXiv preprint arXiv:2410.10934}, 2024.

\bibitem{agentbench}
X.~Liu, H.~Yu, H.~Zhang, et~al.
\newblock {AgentBench}: Evaluating {LLMs} as Agents.
\newblock In {\em Proc.\ ICLR}, 2024.

\bibitem{openagentsafety}
S.~Vijayvargiya, A.~B.~Soni, X.~Zhou, Z.~Z.~Wang, N.~Dziri, G.~Neubig, and
  M.~Sap.
\newblock {OpenAgentSafety}: A Comprehensive Framework for Evaluating
  Real-World {AI} Agent Safety.
\newblock {\em arXiv preprint arXiv:2507.06134}, 2025.

\bibitem{toolemu}
Y.~Ruan, H.~Dong, A.~Wang, S.~Pitis, et~al.
\newblock {ToolEmu}: Identifying Risks of {LLM} Agents with an {LLM}-Emulated
  Sandbox.
\newblock In {\em Proc.\ ICLR}, 2024.

\bibitem{ml_monitoring}
J.~Leest, I.~Gerostathopoulos, P.~Lago, and C.~Raibulet.
\newblock Monitoring and Observability of Machine Learning Systems: Current
  Practices and Gaps.
\newblock {\em arXiv preprint arXiv:2510.24142}, 2025.

\bibitem{aiops_survey}
L.~Zhang, T.~Jia, M.~Jia, Y.~Wu, A.~Liu, Y.~Yang, Z.~Wu, X.~Hu, P.~S.~Yu,
  and Y.~Li.
\newblock A Survey of {AIOps} in the Era of Large Language Models.
\newblock {\em arXiv preprint arXiv:2507.12472}, 2025.

\bibitem{opsagent}
Y.~Luo, J.~Jiang, J.~Feng, L.~Tao, Q.~Zhang, X.~Wen, Y.~Sun, S.~Zhang, and
  D.~Pei.
\newblock From Observability Data to Diagnosis: An Evolving Multi-agent System
  for Incident Management in Cloud Systems.
\newblock {\em arXiv preprint arXiv:2510.24145}, 2025.

\bibitem{otel}
{OpenTelemetry Authors}.
\newblock {OpenTelemetry Specification}, v1.30.
\newblock \url{https://opentelemetry.io/docs/specs/otel/}, 2024.

\bibitem{trace_opt}
C.-A.~Cheng, A.~Nie, and A.~Swaminathan.
\newblock Trace: Unlocking Efficient Optimization of Computational Workflows.
\newblock {\em arXiv preprint arXiv:2406.16218}, 2024.

\bibitem{tracediag}
R.~Ding, C.~Zhang, L.~Wang, Y.~Xu, M.~Ma, W.~Zhang, S.~Qin, S.~Rajmohan,
  Q.~Lin, and D.~Zhang.
\newblock {TraceDiag}: Adaptive, Interpretable, and Efficient Root Cause
  Analysis on Large-Scale Microservice Systems.
\newblock In {\em Proc.\ ACM ESEC/FSE}, 2023.

\bibitem{deeptralog}
C.~Zhang, X.~Peng, C.~Sha, K.~Zhang, Z.~Fu, X.~Wu, Q.~Lin, and D.~Zhang.
\newblock {DeepTraLog}: Trace-Log Combined Microservice Anomaly Detection
  through Graph-based Deep Learning.
\newblock In {\em Proc.\ ACM ICSE}, 2022.

\bibitem{sieve_sampling}
D.~Batavia, I.~Weber, et~al.
\newblock {Sieve}: Attention-based Sampling of End-to-End Trace Data in
  Distributed Microservice Systems.
\newblock In {\em Proc.\ IEEE ICWS}, 2023.

\bibitem{mulan}
L.~Zheng, Z.~Chen, J.~He, and H.~Chen.
\newblock {MULAN}: Multi-modal Causal Structure Learning for Root Cause
  Analysis.
\newblock In {\em Proc.\ NeurIPS}, 2024.

\bibitem{rcagent}
Z.~Wang, Z.~Liu, Y.~Zhang, A.~Zhong, J.~Wang, F.~Yin, L.~Fan, L.~Wu, and
  Q.~Wen.
\newblock {RCAgent}: Cloud Root Cause Analysis by Autonomous Agents with
  Tool-Augmented {LLMs}.
\newblock {\em arXiv preprint arXiv:2310.16340}, 2023.

\bibitem{aiopsbench}
Y.~Liu, C.~Pei, L.~Xu, et~al.
\newblock {OpsEval}: A Comprehensive {IT} Operations Benchmark Suite for
  {LLMs}.
\newblock {\em arXiv preprint arXiv:2310.07637}, 2023.

\bibitem{aioslab}
{Microsoft Research}.
\newblock {AIOpsLab}: A Holistic Framework to Evaluate {AI} Agents for Enabling
  Autonomous Cloud Operations.
\newblock {\em arXiv preprint arXiv:2407.12165}, 2024.

\end{thebibliography}

\end{document}
